% =========================================================
% Proof Stub: Existence of Subsequential Limits
% Source: volume-iii/analysis/sequences/notes/sequences/notes-subsequences.tex
% Mode: standard
% =========================================================

\clearpage
\phantomsection
\hypertarget{proof-RA-ABB-C-subseq-limits-exist}{}

\subsubsection[Existence of Subsequential Limits]{Proof --- Existence of Subsequential Limits}

\bigskip

\noindent
\textbf{Source.}
Abbott, \emph{Understanding Analysis}; see \texttt{notes-subsequences.tex}.

\vspace{0.75em}

\noindent
\textbf{Goal.}
Every bounded sequence admits at least one subsequential limit.

\vspace{0.75em}

\noindent
\textbf{Logical form.}
\[
  \text{bounded } (a_n) \Rightarrow \mathcal{L}(a_n) \ne \varnothing.
\]

\vspace{0.75em}

\noindent
\textbf{Proof strategy.}
By Bolzano--Weierstrass, every bounded sequence has a convergent subsequence, and the limit of that subsequence is a subsequential limit.

\vspace{0.75em}

\noindent
\textbf{Proof.}
\begin{proof}
\Claim Every bounded sequence admits at least one subsequential limit.

\Given 

\Goal 

\Strategy By Bolzano--Weierstrass, every bounded sequence has a convergent subsequence, and the limit of that subsequence is a subsequential limit.

\medskip

% --- apply Bolzano–Weierstrass: get convergent subsequence (a_{n_k}) → L ---
% --- L ∈ L(a_n) by definition of subsequential limit ---
% --- therefore L(a_n) ≠ ∅ ---

\end{proof}

\begin{remark*}[Proof shape]
One-liner corollary: BW gives a convergent subsequence, whose limit is a subsequential limit.
\end{remark*}

\begin{remark*}[Dependencies]
\begin{itemize}
  \item Bolzano--Weierstrass Theorem.
  \item Definition of subsequential limit.
\end{itemize}
\end{remark*}
